\section{\systemname{} Deployment}
The magazine paper~\cite{bao2026llmhric} established the \systemname{} tier decomposition---a non-real-time LLM rApp that reasons over slow, semantically rich summaries, and a near-real-time deep reinforcement learning xApp that acts on fast RAN state---using a transmit-power allocation use case in an IAB network, with a Llama-3.1-8B rApp and a DDPG xApp. This appendix documents what had to change to make that decomposition control \emph{PRBs across two S-NSSAIs} on a live OAI/FlexRIC stack. The tier decomposition, the guidance-conditioned action space, and the asynchronous training discipline are inherited unchanged. Everything else is new: the actuated resource is the per-slice downlink PRB budget introduced in Section~III, the guidance model is a locally served Gemma checkpoint~\cite{gemma2024} rather than a hosted Llama, and the entire observation path is rebuilt on FlexRIC~\cite{flexric2026} E2 service models because no equivalent of the simulated IAB state exists in OAI.

Seven independent clocks coexist in the deployment (Fig.~\ref{fig:timing}, Table~\ref{tab:apps}): RAN slot processing in the gNB; 10~ms E2 service-model reporting; 50~ms monitor summarization, which also sets the nominal width of a derived state window; 100~ms near-RT control; 1000~ms E2SM-KPM audit collection; 10~s rApp language inference; and a free-running asynchronous learner. Fig.~\ref{fig:timing} draws the four clocks that lie on the closed control loop---summarization, control, inference, and learning---and omits the slot, E2-reporting and audit clocks, which are not on it. None of the six RIC-side clocks gates the gNB slot loop. All of them meet in a single SQLite file (\texttt{db\_path} = \texttt{/tmp/llm\_hric/llm\_hric.sqlite3}) opened concurrently by six processes.

\subsection{Measurement and state database}
Three writers populate the measurement tables, and they are separate operating-system processes rather than threads of one collector. The first is the service-model monitor xApp, \texttt{xapp\_mac\_rlc\_pdcp\_gtp\_moni.py}, a modified version of the FlexRIC Python example; it delegates all persistence to the library class \texttt{RicSmDbWriter} in \texttt{ric\_sm\_db\_writer.py}, which performs the windowing and derivation described below under a re-entrant lock because the four service-model callbacks execute on separate SDK threads. The second is the E2SM-KPM collector \texttt{xapp\_kpm\_moni.c}, extended with a SQLite writer and linked against \texttt{-lsqlite3}. The third is the near-RT controller, whose tables are described in the following subsections. The launcher brings up the whole stack with the form used by the repository README, \texttt{UE\_MODE=multi UE\_COUNT=5 START\_LLM\_HRIC=1 START\_GRAFANA=1 DDPG\_APPLY=1 ./run\_e2e\_rfsim.sh start}. The two monitors need no flag of their own---\texttt{START\_KPM\_MONITOR} already defaults to~1---whereas \texttt{START\_LLM\_HRIC} defaults to~0, so omitting it starts the observation path but neither the A1-like server, the rApp, the near-RT controller nor the actuator, and leaves no closed loop. Independently of that flag, the KPM writer stays disabled unless the environment variable \texttt{LLM\_HRIC\_DB\_PATH} is set, which the campaign runner always does.

The monitor xApp subscribes to the FlexRIC MAC, RLC, PDCP and GTP service models at a 10~ms report interval (\texttt{monitor.subscription\_interval\_ms}, restricted by the SDK to $\{1,2,5,10\}$~ms) and admits one indication per model per 50~ms sampling window, dropping indications whose statistics list is empty. Each admitted indication is stored in a layer-specific raw table---\texttt{mac\_ue\_stats\_raw}, \texttt{rlc\_rb\_stats\_raw}, \texttt{pdcp\_rb\_stats\_raw}, \texttt{gtp\_tunnel\_stats\_raw}---carrying both the E2 source timestamp \texttt{ts\_ms}, normalised to milliseconds from the agent's \texttt{tstamp} field, and the RIC receive timestamp \texttt{recv\_ts\_ms}. The 50~ms throttle is a RIC-side observation decimation only: it does not alter the gNB slot, scheduler, PHY or MAC execution period, and the E2 subscription interval likewise controls indication reporting only.

A windowing stage keyed on MAC indication timestamps derives per-UE throughput and slice summaries. Windows are contiguous and half-open, $(t_{k-1},t_k]$, where $t_k$ is the source timestamp of the $k$-th admitted MAC indication; the configured 50~ms sampling period (\texttt{monitor.summary\_period\_ms}) sets the nominal window width and bootstraps the first window only, so realised widths are 50--60~ms rather than a fixed wall-clock grid. Each MAC indication with $t_k>t_{k-1}$ advances the window; an empty UE list is treated as a valid topology state and still advances it, so consumers observe UE departure instead of a stale last-known state. Late-arriving RLC or PDCP indications bearing the same millisecond timestamp cause the window to be re-derived in place: the derived rows at that timestamp are deleted and recomputed so that bearer counters replace the initial MAC fallback. GTP indications are stored but never trigger a flush. The slice rows of every window are pre-seeded from the configured \texttt{slice\_catalog}, so both S-NSSAIs always receive a summary row even when they carry no UE.

Downlink and uplink throughput are computed from counter differences, $R = 8\,\Delta b / (1000\,\Delta t)$~Mbit/s for $\Delta b$ bytes over $\Delta t$~ms, and selected from PDCP, RLC and MAC counters in descending priority: the first strictly positive rate wins, otherwise the first defined rate, with a final fallback of zero. Inside the MAC tier the aggregate MAC-SDU byte counter takes precedence over the transport-block counter, so the recorded provenance source is one of \texttt{pdcp}, \texttt{rlc}, \texttt{mac} or \texttt{zero} and MAC-TBS is not a separately distinguishable source. Bearer-tier rates are summed over radio bearers; each UE here has a single PDU session and therefore a single DRB. A negative counter delta invalidates that tier and additionally raises a \texttt{counter\_reset} provenance flag, which is recorded rather than used to suppress the sample. One deployment fact must be stated plainly: in the monolithic gNB used here the PDCP and GTP service models key their statistics by RRC UE identity rather than by C-RNTI, because \texttt{nr-softmodem} links the combined DU/CU-CP/CU-UP E2 RAN function and the PDCP and GTP handlers populate the identity field from \texttt{rrc\_ue\_id}, a small linear counter. Those rows can never be joined to a UE, so \texttt{pdcp\_rb\_stats\_raw} and \texttt{gtp\_tunnel\_stats\_raw} are populated but unresolvable, and the effective precedence chain reduces to RLC, then MAC-SDU, then MAC-TBS, then zero. This is a property of the stock OAI E2 agent, not of the controller.

MAC counters supply per-window PRB usage and downlink BLER. PRB usage is the delta of the cumulative allocated-RB counter \texttt{dl\_aggr\_prb}, i.e.\ a count of resource blocks in the window and not a share; the slice value is the sum of the per-UE deltas, and the PRB \emph{share} used by the controller and by the rApp prompt is formed downstream by normalising a slice's PRB delta by the cell total at the same timestamp. Slice BLER is the unweighted mean over UEs reporting a value. RLC transmit-buffer occupancy is treated as a gauge rather than a counter: the last sample per bearer in the window is summed, and the slice value is suppressed unless every active UE contributed a sample. Wideband CQI and BLER share a single validity mask requiring $0\leq\mathrm{CQI}\leq15$ and $0\leq\mathrm{BLER}\leq1$ for \emph{every} active UE in the slice; when the mask is not asserted the controller zeroes both features. We note explicitly that in the reported rfsimulator pilot the gNB reports a constant zero wideband CQI and an effectively zero BLER while the mask remains asserted, because zero is inside both admissible ranges. The mask detects absent or out-of-range reports, not degenerate constant ones, and the three channel-related state features consequently carried no information in these runs.

RNTI-to-UE/S-NSSAI bindings are discovered at runtime from the per-UE \texttt{nrUE$i$.log} files, which are the only source that ties a stable UE configuration to its current C-RNTI; the writer matches the stock OAI print \texttt{UE 0 RNTI \%04x stats} in the last 8~MiB of each log and takes the most recent occurrence. The gNB log is used solely as a consistency check against the stock \texttt{UE RNTI \%04x CU-UE-ID} print, and the map is accepted only if every configured UE yields a match, all RNTIs are distinct, and the discovered set is a subset of the RNTIs the gNB reports; otherwise the whole map is refused and the writer falls back to the static list \texttt{monitor.rnti\_slice\_map}. The design reason for preferring the UE-side log is recorded in the code: CU-UE-ID reflects attach order, not catalogue identity, so an attach-order mapping would silently mislabel slices after a UE restart. Indications whose RNTI cannot be resolved are still written to the raw tables with null \texttt{ue\_id}, \texttt{plmn}, \texttt{sst} and \texttt{sd}, and are excluded from every derived table both by an SQL predicate and by an explicit guard, so they remain auditable but are never attributed to a slice. Two deployment caveats follow. First, the gNB cross-check relies on OAI's periodic MAC statistics print, which fires only every 128 frames and which OAI permanently disables once the UE count exceeds \texttt{stats\_max\_ue} (default 8, not overridden by the deployed gNB configuration); the check is therefore live at five UEs and silently degrades to no check above eight. Second, because the PDCP and GTP models report a permanently unresolvable identity, the writer's once-per-second rediscovery path fires for the entire run, re-reading up to six log tails per second.

Every derived UE row is accompanied by a provenance row in \texttt{ue\_metric\_provenance} recording the window bounds, the computation timestamp, the selected downlink and uplink sources, the number of contributing raw samples, the \texttt{counter\_reset} flag and a mapping-validity flag. The near-RT arm consumes these rows directly as a data-quality gate on training transitions, so measurement provenance is a first-class part of the control loop rather than an offline annotation.

A parallel E2SM-KPM collector subscribes to every report style the E2 node advertises and writes a long-format table \texttt{kpm\_measurements\_raw} with columns for scope, UE identity type, measurement name, value type, value, 3GPP unit, labels and a reliability flag derived from the indication's incomplete flag. For the monolithic gNB this is per-UE report style~4 filtered on S-NSSAI (\texttt{DRB.PdcpSduVolumeDL/UL}, \texttt{DRB.RlcSduDelayDl}, \texttt{DRB.UEThpDl/Ul}, \texttt{RRU.PrbTotDl/Ul}) plus cell-scope report style~1 (\texttt{CARR.PDSCHMCSDist}), at an event-trigger and granularity period of 1000~ms that is hard-coded in the collector rather than configured. The S-NSSAI matching condition is a single-octet equality test on SST~1, which both deployed S-NSSAIs satisfy, so all five UEs are reported. No component of the control path reads this table: the control state is assembled exclusively from \texttt{network\_state}, \texttt{applied\_prb\_policy} and \texttt{ue\_metric\_provenance}, and the only reader of \texttt{kpm\_measurements\_raw} anywhere in the source tree is the post-hoc analysis script. KPM is therefore an independent, standards-aligned audit path, and the database must not be described as KPM-only. Two audit-path limitations should be recorded: \texttt{RRU.PrbTotDl/Ul} is reported by the agent as an integer percentage of the node PRB aggregate rather than as an absolute count, and although the action definition requests $8\times3\times32$ distribution bins for \texttt{CARR.PDSCHMCSDist}, the collector's read loop iterates over measurement names rather than label bins, so only the first bin is persisted and it is stored without its label.

Table~\ref{tab:db-tables} groups the principal tables of the schema, which declares 23 tables and 30 indexes and is created idempotently by every process that opens the file. Because a run writes on the order of $700$~rows/s for roughly 90~minutes into a single WAL-mode SQLite file shared by six processes, indexing is a deployment requirement rather than an optimisation. Composite indexes \texttt{(sst,sd,ts\_ms)} on \texttt{network\_state} and \texttt{(sst,sd,applied,ts\_ms)} on \texttt{applied\_prb\_policy} were added specifically to keep the 100~ms controller tick from drifting: without them SQLite scans a table growing at 40~rows/s for the whole run on every tick, and tick jitter grows without bound. Similarly, \texttt{(rnti,ts\_ms)} and \texttt{(rnti,rbid,ts\_ms)} indexes on the raw tables keep the windowing stage's previous-sample lookups constant-time. Connections are opened with \texttt{journal\_mode=WAL}, \texttt{synchronous=NORMAL} and \texttt{busy\_timeout=5000}, the WAL pragma is retried for 5~s because several processes may race to set it, and schema evolution uses idempotent \texttt{ALTER TABLE} helpers that tolerate a competing process having added the same column. Each pilot run's database is 1.0--1.3~GB for an approximately 88~minute run.

\begin{table*}[t]
\centering
\caption{Principal tables of the shared SQLite state database. Design rates assume five UEs, one DRB each, two slices, and 20 admitted samples per second per service model; they are properties of the sampling design, not measured counts.}
\label{tab:db-tables}
\footnotesize
\begin{tabular}{llp{0.46\textwidth}l}
\toprule
Category & Table & Purpose & Design rate\\
\midrule
\multirow{5}{*}{Raw}
 & \texttt{mac\_ue\_stats\_raw} & Per-UE MAC counters per admitted indication: TBS and SDU bytes, aggregate and scheduled RBs, BLER, SNR, wideband CQI, MCS, BSR, PHR & 100~s$^{-1}$\\
 & \texttt{rlc\_rb\_stats\_raw} & Per-DRB RLC transmitted/received SDU bytes and transmit-buffer occupancy & 100~s$^{-1}$\\
 & \texttt{pdcp\_rb\_stats\_raw} & Per-DRB PDCP SDU bytes; keyed by RRC UE identity in this gNB and therefore unjoinable & 100~s$^{-1}$\\
 & \texttt{gtp\_tunnel\_stats\_raw} & Per-QFI GTP tunnel endpoints; unjoinable for the same reason & 100~s$^{-1}$\\
 & \texttt{kpm\_measurements\_raw} & Long-format E2SM-KPM records with 3GPP unit, labels and reliability flag & 36~s$^{-1}$\\
\midrule
\multirow{3}{*}{Derived}
 & \texttt{network\_state} & Per-slice window summary: UE count, DL/UL throughput, PRB usage, BLER, RLC buffer, wideband CQI, channel-valid mask & 40~s$^{-1}$\\
 & \texttt{ue\_slice\_throughput} & Per-UE DL/UL throughput at the window end & 100~s$^{-1}$\\
 & \texttt{ue\_metric\_provenance} & Per-UE window bounds, selected throughput sources, contributing sample count, counter-reset and mapping-validity flags & 100~s$^{-1}$\\
\midrule
\multirow{6}{*}{Policy}
 & \texttt{intent\_inputs} & Versioned operator intents, one active row per intent identifier & event\\
 & \texttt{llm\_guidance} & One audit row per accepted inference: model identifier, prompt SHA-256, parsed guidance & $0.1$~s$^{-1}$\\
 & \texttt{a1\_policies} & Versioned A1-like policy payloads, exactly one active row per policy identifier & event\\
 & \texttt{a1\_policy\_summary} & Flattened view of the active policy for operator tooling & event\\
 & \texttt{ddpg\_actions} & Issued near-RT actions with state, raw and scaled reward & control tick\\
 & \texttt{applied\_prb\_policy} & Per-slice dedicated PRB ratio acknowledged by the RAN & control tick\\
\midrule
\multirow{5}{*}{Audit}
 & \texttt{experiment\_runs} & Run identity: arm, seed, scenario, intent, status, configuration, git revision & per run\\
 & \texttt{experiment\_steps} & Per-step transition record with reward components, control latency, effect coverage, validity & control tick\\
 & \texttt{traffic\_events} & Planned and applied per-UE offered-load changes with trace hash & per phase\\
 & \texttt{controller\_ticks} & Scheduled and actual wake times, jitter, state freshness, skip reason, policy version & control tick\\
 & \texttt{channel\_events} & Planned and applied rfsimulator channel perturbations (disabled in the reported runs) & event\\
\midrule
\multirow{4}{*}{Learner}
 & \texttt{ddpg\_replay\_transitions} & Raw replay rows with contexts, policy version, effect coverage and validity flag & control tick\\
 & \texttt{ddpg\_runtime\_state} & Learner cursor, valid replay count, pending updates, serving version, heartbeat & continuous\\
 & \texttt{ddpg\_actor\_versions} & Candidate Actor snapshots with SHA-256, acceptance decision and reason & per candidate\\
 & \texttt{ddpg\_learner\_updates} & Per-train-step critic and actor losses, TD error, predicted and target $Q$ & per train step\\
\bottomrule
\end{tabular}
\end{table*}

\subsection{LLM rApp and A1-like policy}
The rApp is a single Python process, \texttt{llm\_guidance\_service.py}, started by the launcher one second after the A1-like policy server and enabled for \emph{all three} experimental arms so that the arms differ only in how the near-RT controller consumes guidance. Its schedule is a 10~s timer (\texttt{period\_ms.llm} = 10000) overlaid on a 200~ms intent poll: on each poll it reads the active intent and runs inference if either the periodic deadline has passed or the intent version has changed. A newly posted intent therefore takes effect within one poll interval and off the nominal cadence, without a restart. If an inference overruns the period the scheduler advances to the next 10~s boundary rather than queuing missed cycles.

Each cycle reads the three most recent ordered, non-overlapping 10~s windows, anchored on the newest slice summary in \texttt{network\_state} rather than on wall clock: the query selects rows in $(\,t_{\max}-30000,\,t_{\max}]$ and buckets them into windows indexed $0$ (oldest) to $2$. At the 50~ms monitor cadence each window contains about 200 samples per slice. For each slice and window the prompt carries exactly six fields: a compact S-NSSAI key of the form \texttt{sst:sd} with the \texttt{0x} prefix removed (\snssaiA, \snssaiB), the end-of-window UE count, the window-mean downlink throughput in Mbit/s, the window-mean \emph{inter-slice} PRB share, the window-mean RLC downlink buffer occupancy in bytes, and a per-slice availability flag; each window additionally carries its index, start and end timestamps, and its own availability flag. The PRB share is computed per timestamp as $\mathrm{prb\_used}_i/\sum_j\mathrm{prb\_used}_j$ and then averaged, so it sums to one across slices and is an inter-slice share rather than a cell-utilisation figure. A window is marked available only if every catalogue slice in it has a defined throughput mean, a defined buffer mean, at least one share sample, and a sample span covering at least 80\% of the nominal window length. BLER, wideband CQI, latency and all uplink counters are present in \texttt{network\_state} but are deliberately withheld from the prompt.

The provider layer is pluggable across three back ends: a deterministic mock, an OpenAI-compatible HTTP endpoint, and local HuggingFace \texttt{transformers}. The deployed configuration, summarised in Table~\ref{tab:llm-params}, selects \texttt{transformers} with \texttt{google/gemma-4-E2B-it}~\cite{gemma2024}, loaded entirely from local files onto a single GPU under bitsandbytes NF4 4-bit quantisation with double quantisation and float16 compute. Decoding is greedy: a configured temperature of $0$ sets \texttt{do\_sample=false} and removes the temperature argument, and the client never sets \texttt{top\_p} or \texttt{top\_k}, so the checkpoint's sampling defaults are inert. Generation is bounded twice, by \texttt{max\_new\_tokens}=256 and by a 30~s wall-clock generation cap. The client is cached on the triple (provider, model, base URL), so the checkpoint is loaded once per rApp process. Prompts are wrapped with the processor's chat template using the system message \texttt{"Return only JSON guidance for an O-RAN near-RT controller."} and the generated continuation is decoded after stripping the prompt tokens. One measurement gap must be stated plainly, because it bears directly on the architectural premise this appendix inherits. The central argument of the magazine paper is that LLM inference is too slow to sit on the near-RT scheduling path, yet rApp inference latency was \emph{not} instrumented in this deployment: the guidance service times a cycle only to advance its own next deadline, it never persists or logs that duration, the \texttt{llm\_guidance} table has no latency column, and no surviving derived metric carries one. The 10~s rApp period reported here is therefore a design choice, not a measured lower bound on the guidance model's response time, and the only latency-relevant constants in the system are that period and the 30~s generation cap. Instrumenting per-cycle generation latency is a prerequisite for any quantitative restatement of the tier-separation argument on this stack. Relative to the magazine deployment, this replaces a Llama-3.1-8B guidance model for power allocation with a quantised, locally served Gemma checkpoint for PRB guidance; the interface contract between the tiers is unchanged.

\begin{table}[t]
\centering
\caption{Deployed rApp inference configuration (\texttt{config.yaml}, key \texttt{llm}). The file is parsed as JSON despite its extension.}
\label{tab:llm-params}
\footnotesize
\begin{tabular}{ll}
\toprule
Key & Deployed value\\
\midrule
\texttt{provider} & \texttt{transformers}\\
\texttt{model}, \texttt{processor\_model} & \texttt{google/gemma-4-E2B-it}\\
\texttt{model\_class} & \texttt{AutoModelForCausalLM}\\
\texttt{device\_map} & single GPU (device 0)\\
\texttt{torch\_dtype} / \texttt{quantization} & \texttt{float16} / \texttt{4bit} (NF4, double)\\
\texttt{max\_new\_tokens} & 256\\
\texttt{temperature} & $0.0$ (greedy, \texttt{do\_sample=false})\\
\texttt{max\_time\_s} & 30 (generation wall-clock cap)\\
\texttt{local\_files\_only} & \texttt{true} (fully offline load)\\
\texttt{trust\_remote\_code} & \texttt{false}\\
\texttt{allow\_fallback} & \texttt{false}\\
\texttt{require\_real\_model} & \texttt{true}\\
\texttt{intent\_poll\_ms} & 200\\
\texttt{state\_window\_ms} / \texttt{\_count} & 10000 / 3\\
\texttt{guidance\_bound\_tolerance\_prb} & 10\\
\bottomrule
\end{tabular}
\end{table}

Fig.~\ref{fig:prompt} reproduces the prompt template and its substituted payloads. The prompt is emitted as a single line with no newlines and interpolates exactly five payloads: the operator intent verbatim, the ordered slice-key list in catalogue order, the three-window compact state, the calibrated policy-constraint object, and a shape example that is generated at run time from the catalogue rather than hard-coded. The policy-constraint object is the material addition relative to the magazine prompt: it hands the model the priority and protected S-NSSAI identities, the calibrated cell capacity, the SLA floor, the minimum priority PRB ratio, the feasible PRB bounds, the calibration diagnostics used to derive them, \emph{and the calibrated preferred ratio itself}. The model's role is therefore bounded refinement inside a pre-certified feasible region rather than unconstrained allocation. Two details are worth recording for exact reproduction: the slice-key list is in catalogue order while the compact-state array is in sorted key order, and the \texttt{priority\_min\_gap\_prb} configuration value is accepted by the prompt builder but never appears in the prompt body.

\begin{figure}[t]
\centering
\scriptsize
\begin{verbatim}
You are an O-RAN rAPP. Return exactly one minified
JSON object and no markdown. Choose integer DL PRB
ratios for the listed slice keys. The ratios must be
0..100 and sum exactly 100. Choose a preferred ratio
that protects the throughput floor while maximizing
the priority slice. Intent: {INTENT}. Slice keys:
{KEYS}. Three ordered 10-second state windows JSON
(oldest first): {STATE}. Calibrated policy
constraints JSON: {CONSTRAINTS}. Required JSON shape
example: {EXAMPLE}. Return only control_type and
action using that exact compact shape. Replace the
example ratios with your chosen action. Use numbers,
not strings. Keep action.reason to at most four
words.

{INTENT} = Maximize total DL throughput, prioritize
  slice 1:ffffff, and keep aggregate DL throughput
  of slice 1:123456 at or above 24.31 Mbps.
{KEYS}   = 1:ffffff,1:123456
{STATE}  = [{"i":0,"start_ms":...,"end_ms":...,
   "available":true,"slices":[
     {"k":"1:123456","ue":2,"dl":33.46,
      "prb_share":0.339,"rlc_buf_bytes":24986.3,
      "available":true},
     {"k":"1:ffffff","ue":3,"dl":61.96,
      "prb_share":0.661,"rlc_buf_bytes":40861.1,
      "available":true}]},{"i":1,...},{"i":2,...}]
{CONSTRAINTS} = {"sla_floor_mbps":24.3131,
   "cell_capacity_mbps":99.2412,
   "priority_min_ratio":55,
   "priority_slice":"1:ffffff",
   "protected_slice":"1:123456",
   "preferred_ratio":{"1:ffffff":65,"1:123456":35},
   "bounds":{"1:ffffff":{"min_prb":55,"max_prb":75},
             "1:123456":{"min_prb":25,"max_prb":45}},
   "protected_p10_at_target":27.0146,
   "safe_sweep_ratios":[55,65,75]}
{EXAMPLE} = {"control_type":"prb","action":
   {"prb_ratio":{"1:ffffff":50,"1:123456":50},
    "confidence":0.8,"reason":"short"}}
\end{verbatim}
\caption{rApp prompt template and its five interpolated payloads. The template text is verbatim from \texttt{build\_prompt}; the payload values are illustrative and structurally exact. Prompts cannot be replayed from a run database because only the SHA-256 of the prompt is stored.}
\label{fig:prompt}
\end{figure}

The model is asked for a single minified JSON object of the form \texttt{\{"control\_type":"prb","action":\{"prb\_ratio":\{\dots\},"confidence":$c$,"reason":$s$\}\}}, that is, integer target PRB ratios per slice key, a confidence, and a rationale of at most four words. Bounds are \emph{not} requested from the model. The extractor is tolerant by design: it strips markdown fences, falls back to the first brace-delimited object found in the text, and closes containers left open by the 256-token generation cap by walking the string with quote, escape and bracket-stack tracking. Validation then rejects a non-object ratio map, an unknown slice key, a non-numeric or out-of-range value, and a ratio sum exceeding 100; a sum \emph{below} 100 is accepted and completed by the projection step, and catalogue keys absent from the response default to zero. Confidence is coerced to $[0,1]$ with a default of $1.0$. The rApp then derives a $\pm10$~PRB envelope around the returned ratios, intersects it with any bounds the model volunteered, and---whenever calibrated policy constraints are present, which is always in the campaign---discards both and substitutes the calibrated feasible bounds. The returned ratios are finally projected into that region: values are clipped into $[l_k,u_k]$, and the residual to a total of exactly 100 is distributed one PRB at a time to the slice with the largest remaining deficit (or surplus), with ties broken toward the earlier catalogue index; a region that cannot reach a total of 100 raises an error rather than silently renormalising. The raw pre-projection ratios are retained in the published payload for audit.

The rApp fails safe by holding, not by inventing. There is no heuristic default-policy generator in the live path. If provider construction fails---missing weights, absent bitsandbytes, or out-of-memory---the mock client is substituted only when \texttt{allow\_fallback} is true \emph{and} \texttt{require\_real\_model} is false; the campaign sets the opposite, so the rApp refuses to start rather than silently degrade. This safety property has a reproducibility cost that must be recorded: with \texttt{local\_files\_only} the loader resolves \texttt{google/gemma-4-E2B-it} strictly against the local HuggingFace cache, and that identifier does not correspond to a resolvable public repository, so a third party cannot presently obtain the exact weights this pilot used and cannot start the rApp at all without substituting a checkpoint of their own.
% TODO(authors): record the true upstream checkpoint identifier and its exact
% revision/commit hash for the weights actually served in the pilot, and update
% both this sentence and Table~\ref{tab:llm-params} accordingly. If the model output cannot be parsed as JSON even after fence stripping, object extraction and truncation repair, the cycle aborts, the service logs that it is preserving the last active A1 policy, and no new policy version is written, so the previously activated policy remains in force. If the output parses but the action is semantically invalid---unknown S-NSSAI, ratios summing above 100, malformed bounds, or a projection that cannot reach 100---the rApp emits an audit line and returns the currently active policy unchanged; only when no policy has ever been activated does it propagate the error. On the near-RT side, a stored policy that cannot be re-parsed degrades to the stored preferred ratio at zero confidence with unconstrained bounds, and the complete absence of an active policy degrades to a uniform split across the catalogue. Finally, policy lifetime is recorded but not enforced: every intent and policy payload carries \texttt{valid\_for\_ms}, but nothing reads it, so the last activated policy is authoritative until superseded.

Intents enter the system by one of two paths. At boot the launcher passes a seed intent, which is written only if no active intent exists; at run time the campaign runner issues \texttt{POST /a1-p/intents/slice-prb-intent} with the intent text, the slice scope, and the calibrated policy constraints, and the 200~ms poll picks up the version bump. Crucially, the safety-relevant semantics of the intent---which S-NSSAI is prioritised, which is protected, and the numeric SLA floor---are \emph{not} parsed by the LLM. They are taken from the structured calibration constraints when present and otherwise lifted by fixed regular expressions (a \texttt{prioriti[sz]e \dots} pattern matched against each catalogue key, and an \texttt{(above$\vert$at or above$\vert$>=) $x$ Mbps} pattern for the floor), with the protected slice defaulting to the complement of the priority slice in a two-slice catalogue. The near-RT controller re-derives the same quantities independently with the same precedence and the same patterns. The LLM therefore influences only the PRB split inside an already-certified feasible region; it cannot redefine what the operator asked for.

The published policy is a versioned document. The rApp first writes an audit row into \texttt{llm\_guidance} recording the timestamp, policy identifier, intent, the parsed guidance object, the model identifier and the SHA-256 of the prompt, then POSTs the policy payload to the A1-like server. Because the prompt is reduced to a digest and the raw completion is discarded once parsed, no rApp inference from this pilot can be replayed or audited from the artifact: the digest establishes only that two cycles used byte-identical prompts, and neither the state payload that produced a given guidance object nor the text the model actually returned is recoverable. Archiving the full prompt and the full completion alongside the digest is a requirement for reproducibility of the guidance path, and is not satisfied here. That payload preserves the untouched model output under \texttt{guidance.action} and grafts on four derived keys: \texttt{preferred\_ratio} (the raw pre-projection ratios), \texttt{a1\_target\_ratio} (the projected ratios), \texttt{bounds} (the calibrated feasible region), and \texttt{policy\_context} (the resolved priority and protected slices, floor, minimum priority ratio and capacity). Activation is a single SQLite transaction: the next version number is taken as one above the current maximum for that policy identifier, all existing rows for the identifier are cleared of the active flag, the new row is inserted active, the flattened summary is upserted, and the transaction commits. Under WAL a concurrent near-RT reader therefore observes either the previous or the new active policy, never a torn state.

This interface is \emph{A1-like}: it realises the policy role this prototype needs but is not an O-RAN A1-PMS implementation~\cite{oran_a1}. The deviations are specific. It exposes \texttt{POST}/\texttt{GET} on \texttt{/a1-p/policies/\{id\}} rather than the A1AP \texttt{PUT /A1-P/v2/policytypes/\{typeId\}/policies/\{id\}}, with no interface-version path segment. There is no policy-type registry and therefore no schema-typed policy objects; the body is an ad-hoc structure. There is no \texttt{DELETE}, no \texttt{/status} resource, no policy-enforcement feedback and no notification-destination callback. There is no TLS and no authentication, and the server binds to loopback only. It adds an \texttt{/a1-p/intents/\{id\}} resource that has no A1-P counterpart, since O-RAN places intent handling with the non-RT RIC rather than on A1. Policy lifetime is carried but not enforced. Finally, the consumer path is not HTTP at all: the near-RT controller reads the active policy row directly from the shared SQLite file, and the HTTP surface is producer- and operator-facing only.
